(1) 电子 Hamiltonian 是
\[
H=\frac1{2m}[p_x^2+(p_y+eBx)^2].
\]
因 y 平移不变，可以取 \(\psi(x,y)=\phi(x)e^{ik_yy}\)，于是 \(p_y\to\hbar k_y\)，x 方向有效势为 \(e^2B^2(x-x_c)^2/(2m)\)，其中 \(x_c=-\hbar k_y/(eB)\)。这与质量 m、频率 \(\omega_c=eB/m\) 的谐振子完全等价。

(2) 所有 Landau 能级为
\[
E_n=\hbar\omega_c(n+1/2),\quad n=0,1,2,\ldots,
\qquad \ell_B=\sqrt{\hbar/(eB)}.
\]
波函数的横向特征长度为磁长 \(\ell B\)；能量不含 ky，因此轨道中心可连续变化并产生宏观简并。

(3) 在 y 长度 Ly 的周期边界中 \(k_y=2\pi j/L_y\)。相邻导引中心间距为 \(\Delta x_c=2\pi\hbar/(eBL_y)=2\pi\ell_B^2/L_y\)。在 x 的样品长度 Lx 内允许的导引中心数为
\[
N_\Phi\simeq\frac{L_x}{\Delta x_c}
=\frac{eB L_xL_y}{h}=\frac{eBS}{h}.
\]
这等于穿过样品的磁通除以磁通量子 h/e；忽略边缘取整误差。

(4) 若 \(\nu \) 个自旋极化的 Landau 级均满，电子二维面密度为 \(n_{2D}=\nu N_\Phi/S=\nu eB/h\)。经典电磁漂移在无耗散稳态中满足霍尔电导大小 \(n_{2D}e/B\)，所以
\[
|\sigma_{xy}|=\nu\frac{e^2}{h}.
\]
此结果体现量子化的占据态数与磁通简并；电子的符号决定霍尔电压方向。真实样品中观测宽阔整数平台还需要能隙、杂质局域态及低温等机制，不能仅用无杂质体相态计数证明整个平台宽度。
